\chapter{Hopf fibrations and toric potentials}
\label{chap:hopf-reduction}

Fix $n\geq2$ and use the scalar Hopf manifold $X_n$, projection $\pi:X_n\to\mathbb P^{n-1}$ and forms $\alpha,\zeta,\omega$ of \cref{prop:hopf-bergman-exact}. Write $\sigma$ for the Fubini--Study form normalized by $\int_{\mathbb P^{n-1}}\sigma^{n-1}=1$.

Potentials in the pulled-back closed classes descend to the projective base because the elliptic fibers are compact. Fiber integration identifies the weighted section norms on both spaces. For descended $\omega$-plurisubharmonic potentials, it also computes the Monge--Amp\`ere measure and reduces the equation to the base. There closedness fixes the normalization and gives uniqueness, although unrestricted masses on the Hopf threefold are unbounded. For smooth base potentials invariant under separate coordinate rotations, logarithmic coordinates turn curvature positivity into convexity; strict positivity on the base makes the gradient parametrize the relative interior of a simplex.

The fiber calculations use the identities, for $y\in\mathbb P^{n-1}$,
\[
 \alpha=\pi^*\sigma,\qquad \omega=\alpha+\zeta,\qquad
 \zeta^2=0,\qquad
 \mathrm{d}V=\tfrac12\alpha^{n-1}\wedge\zeta,\qquad
 \int_{\pi^{-1}(y)}\zeta=2.
\]

\section{Weighted sections on the projective base}\label{sec:hopf-projective-sections}

The line bundles $L_k$ from \cref{prop:hopf-bergman-exact}, for integers $k\geq0$, are the pullbacks of $\mathcal O(k)$ with its Fubini--Study metric. Represent a section by a homogeneous polynomial $f$ of degree $k$ in the covering coordinates $z\in\mathbb C^n\setminus\{0\}$. Its squared norm $|f(z)|^2|z|^{-2k}$ is invariant under every nonzero complex scalar. Descent of the weight therefore makes the weighted integrand a function on the base, and the fiber integral identifies the norms and kernels of \cref{sec:weighted-section-spaces}.

\Needspace{18\baselineskip}
\begin{theorem}\label{thm:hopf-weighted-section-descent}\index{Hopf manifold!section descent}
Let $k\geq0$ be an integer and $u\in\PSH(X_n,k\alpha)$. Then $u=\pi^*v$ for a unique $v\in\PSH(\mathbb P^{n-1},k\sigma)$. For every positive integer $m$, pullback identifies the weighted section spaces isometrically; for a section $s$ in the displayed space on the base,
\[
 \begin{split}
 H^0\bigl(\mathbb P^{n-1},\mathcal O(km)\otimes\mathcal I(mv)\bigr)
 &\simeq H^0\bigl(X_n,L_k^m\otimes\mathcal I(mu)\bigr),\\
 \|\pi^*s\|_{X_n,m,u}^2&=\|s\|_{\mathbb P^{n-1},m,v}^2.
 \end{split}
\]
The norm on the base uses $\sigma^{n-1}$. The kernels $K$ and, off their pole sets, $B$ are pullbacks of the corresponding base kernels. Pullback also identifies the bounded weighted graded section algebras.
\end{theorem}
\begin{proof}
Descent of the potential follows from \cref{prop:psh-descent-submersion}. The unweighted section spaces on both manifolds consist of the same homogeneous polynomials of degree $km$, by the proof of \cref{prop:hopf-bergman-exact}. Pointwise norms of corresponding sections agree by the metric identification. For every nonnegative Borel function $F$ on the base, fiber integration gives
\[
 \int_{X_n}(F\circ\pi)\,\mathrm{d}V
 =\frac12\int_{\mathbb P^{n-1}}F\sigma^{n-1}
                         \int_{\pi^{-1}(y)}\zeta
 =\int_{\mathbb P^{n-1}}F\sigma^{n-1}.
\]
Apply this formula to $F=|s|^2\mathrm{e}^{-mv}$, with the value on the pole set immaterial. It proves equality of the possibly infinite integrals. Finiteness characterizes the multiplier-ideal section spaces, giving the asserted isometry. An orthonormal basis therefore pulls back to an orthonormal basis, which proves the kernel identities. Finally, a nonnegative function on the base is bounded if and only if its pullback is bounded, proving the statement about bounded weighted sections and their multiplication.
\end{proof}

\section{Descended measures and equations}

The form $\zeta$ supplies the fiber direction missing from the pulled-back curvature. Its square vanishes, so the Hermitian top power of a descended potential has a single mixed term. Integrating this term determines the mass and reduces the equation within this class to the base.

If $\nu$ is a positive Borel measure on $\mathbb P^{n-1}$, for nonnegative Borel test functions $F$ on $X_n$ we write $\pi^*\nu\wedge\zeta$ for the measure defined by
\[
 \int_{X_n}F\,\pi^*\nu\wedge\zeta
 =\int_{\mathbb P^{n-1}}
        \left(\int_{\pi^{-1}(y)}F\zeta\right)\mathrm{d}\nu(y).
\]
For a smooth top-degree form $\nu$, this agrees with its pullback wedged with $\zeta$.

\Needspace{18\baselineskip}
\begin{theorem}\label{thm:hopf-descended-hermitian-ma}\index{Hopf manifold!descended equation}
For a quasi-plurisubharmonic function $v$ on $\mathbb P^{n-1}$, the pullback $u=\pi^*v$ belongs to $\PSH(X_n,\omega)$ if and only if $v\in\PSH(\mathbb P^{n-1},\sigma)$. In this case
\[
 \mathrm{MA}^{\mathrm{ext}}_\omega(u)
 =n\,\pi^*\mathrm{MA}^{\mathrm{ext}}_\sigma(v)\wedge\zeta.
\]
Consequently
\[
 \int_{X_n}\mathrm{MA}^{\mathrm{ext}}_\omega(u)
 =2n\int_{\mathbb P^{n-1}}\mathrm{MA}^{\mathrm{ext}}_\sigma(v)
 \leq2n.
\]
Equality holds for bounded $v$. For an integer $k\geq0$ and $v\in\PSH(\mathbb P^{n-1},k\sigma)$,
\[
 \mathrm{MA}^{\mathrm{ext}}_{k\alpha}(\pi^*v)=0
\]
as an $n$-dimensional measure on $X_n$.
\end{theorem}
\begin{proof}
If $v$ is $\sigma$-plurisubharmonic, then
\[
 \omega+\ddc u=\pi^*(\sigma+\ddc v)+\zeta\geq0.
\]
Conversely, if the left side is positive, wedge it with the smooth semipositive form $\zeta$ and push forward. Since $\zeta^2=0$ and its fiber integral is two,
\[
 \pi_*\bigl((\omega+\ddc u)\wedge\zeta\bigr)
 =2(\sigma+\ddc v)\geq0.
\]
This identity of currents follows by testing against a smooth form on the base and integrating in the fiber variables. It proves the converse.

For the measure formula, first take bounded $v\in\PSH(\mathbb P^{n-1},\sigma)$. By \cref{lem:bounded-products-submersion}, its bounded products commute with $\pi^*$. Since $\zeta^2=0$ and the base has dimension $n-1$, expansion gives
\[
 (\omega+\ddc u)^n
 =n\,\pi^*(\sigma+\ddc v)^{n-1}\wedge\zeta.
\]
The definition of the measure on the right agrees with this local product by fiber integration.

For general $v$, normalize $v\leq0$ and use $v_j=\max(v,-j)$. Because $\sigma$ is positive, these are bounded $\sigma$-plurisubharmonic functions. Their pullbacks are exactly $u_j=\max(u,-j)$. Multiply the bounded identity by the common indicator $\mathbf1_{\{u>-j\}}=\pi^*\mathbf1_{\{v>-j\}}$ and pass to the increasing limits defining the extended measures. Fiber integration for nonnegative functions justifies the limit on the right, including at poles.

The mass formula follows because each fiber has $\zeta$-mass two. The base extended mass is at most one: its truncation measure restricted to $\{v>-j\}$ is bounded above in total mass by the full bounded product, whose mass is $\int\sigma^{n-1}=1$ by closed Stokes. For bounded $v$ no part is removed and equality holds.

For the final assertion, take bounded truncations of $v\in\PSH(\mathbb P^{n-1},k\sigma)$. Their pulled-back currents have zero $n$th power by \cref{lem:bounded-products-submersion}. Restricting to the truncation levels and taking the increasing limit gives zero extended measure, also when $k=0$.
\end{proof}

If an $\omega$-plurisubharmonic potential is constant on the fibers, the local-section argument in \cref{prop:psh-descent-submersion} gives a quasi-plurisubharmonic base potential. The theorem therefore bounds its Hermitian mass by $2n$. On the Hopf threefold this bound is six, so the unbounded masses in \cref{chap:hopf} require potentials that vary along the fibers.

For the equation, let $f\geq0$ be in $L^p(\mathbb P^{n-1},\sigma^{n-1})$, $p>1$, with $\int f\sigma^{n-1}>0$. Among potentials $u=\pi^*v$ with bounded $v\in\PSH(\mathbb P^{n-1},\sigma)$,
\begin{equation}\label{eq:hopf-reduced-equation}\index{Monge--Ampere equation@Monge--Amp\`ere equation!Hopf reduction}
 (\omega+\ddc u)^n=c(f\circ\pi)\,\mathrm{d}V,\qquad c>0,
\end{equation}
is equivalent to the base equation
\[
 (\sigma+\ddc v)^{n-1}=\frac{c}{2n}f\sigma^{n-1}.
\]
The base equation implies \eqref{eq:hopf-reduced-equation} by \cref{thm:hopf-descended-hermitian-ma} and the formula for $\mathrm{d}V$. Conversely, pushing \eqref{eq:hopf-reduced-equation} forward gives $2n(\sigma+\ddc v)^{n-1}=cf\sigma^{n-1}$.

\Needspace{5\baselineskip}
Integrating the base equation fixes
\[
 c=\frac{2n}{\int_{\mathbb P^{n-1}}f\sigma^{n-1}}.
\]
Since $\sigma$ is closed, it satisfies the Guan--Li condition. Apply \cref{uni-thm:GL-intro} to $(\mathbb P^{n-1},\sigma)$. It gives a continuous base solution with the scalar just computed, unique up to constants among bounded $\sigma$-plurisubharmonic functions. Consequently \eqref{eq:hopf-reduced-equation} has a continuous solution unique up to constants within the descended class. For smooth positive $f$, \cref{thm:tw-smooth} on the same base gives smooth $v$, with the same scalar fixed by Stokes.

For a real parameter $\lambda$, the same calculation identifies an equation with right side $c\mathrm{e}^{\lambda u}(f\circ\pi)\mathrm{d}V$ with the base equation having right side $(c/(2n))\mathrm{e}^{\lambda v}f\sigma^{n-1}$.

\section{Logarithmic coordinates and the gradient simplex}

Projective scaling leaves one null direction in the logarithmic Hessian. On its quotient, positivity becomes convexity, and boundedness of the potential forces the gradient into a simplex. The logarithmic growth gives every interior gradient value; strict positivity gives uniqueness on the quotient.

Here allow $k>0$ to be real. Let $v$ be a smooth real function on $\mathbb P^{n-1}$ invariant under separate rotations of the homogeneous coordinates, and put $u=\pi^*v$. In the covering coordinates $z_1,\ldots,z_n$, on the locus where all coordinates are nonzero, set $x_i=\log|z_i|^2$ and
\[
 F(x)=k\log\left(\sum_{i=1}^n \mathrm{e}^{x_i}\right)+u(x).
\]
Write $\nabla F=(F_i)$ and $D^2F=(F_{ij})$ for its Euclidean gradient and Hessian, and let $\mathbf1=(1,\ldots,1)$. This invariance makes $u$ a function of $x$, while projective scaling gives, for every $t\in\mathbb R$,
\begin{equation}\label{eq:hopf-convex-translation}
 F(x+t\mathbf1)=F(x)+kt.
\end{equation}

\begin{proposition}\label{prop:hopf-convex-hyperplane}\index{Toric potential}\index{Gradient simplex}
On this locus, positivity of $k\alpha+\ddc u$ is equivalent to convexity of $F$. When $F$ is convex, its gradient lies in
\[
 \Sigma_k=\{p\in\mathbb R^n:p_i\geq0,\ \textstyle\sum_i p_i=k\},
\]
and $D^2F\,\mathbf1=0$. If $k\sigma+\ddc v$ is positive definite on the base, the gradient induces a bijection from $\mathbb R^n/\mathbb R\mathbf1$ onto the relative interior of $\Sigma_k$.
\end{proposition}
\Needspace{8\baselineskip}
\begin{proof}
On the covering space, $k\alpha+\ddc u=\ddc F$. In the coordinates $z_i$, its Hermitian coefficient matrix is a positive constant times
\[
 \left(\frac{F_{ij}}{z_i\overline z_j}\right)_{i,j}.
\]
Congruence by the invertible diagonal matrix with entries $1/z_i$ proves the equivalence with the nonnegative real Hessian of $F$. Differentiating \eqref{eq:hopf-convex-translation} gives $\sum_i F_i=k$ and $D^2F\,\mathbf1=0$.

Boundedness of $v$ on the compact base gives a constant $C$ such that
\[
 \left|F(x)-k\log\sum_i \mathrm{e}^{x_i}\right|\leq C.
\]
Fix $x$ and a coordinate vector $e_i$. The convex function $t\mapsto F(x+te_i)$ is bounded above as $t\to-\infty$, since $n\geq2$. Its derivative cannot be negative: the supporting line at such a point would tend to $+\infty$ as $t\to-\infty$. Thus $F_i\geq0$, proving $\nabla F\in\Sigma_k$.

Assume the base form is positive definite. The Hessian is positive definite transverse to $\mathbf1$, because the differential of projectivization has kernel exactly the complex scalar direction. In particular $F$ is strictly convex on $H=\{x:\sum_i x_i=0\}$. Along each line $x+te_i$, its derivative is strictly increasing and nonnegative everywhere, hence positive. Thus every gradient lies in the relative interior of $\Sigma_k$.

Conversely take $p_i>0$ with $\sum_i p_i=k$. The function $F(x)-p\cdot x$ is translation invariant along $\mathbf1$. On $H$ it tends to $+\infty$ as $|x|\to\infty$. To see this, note that
\[
 k\max_i x_i-p\cdot x
 =\sum_i p_i(\max_j x_j-x_i)
\]
is strictly positive on every nonzero vector of $H$, and has a positive minimum on its unit sphere. The bounded difference between $F$ and $k\log\sum \mathrm{e}^{x_i}$, together with $\log\sum \mathrm{e}^{x_i}\geq\max_i x_i$, gives the claimed coercivity. There is thus a unique minimizer on $H$. Its critical-point equation says that $\nabla F-p$ is a multiple of $\mathbf1$; both vectors have coordinate sum $k$, so this multiple is zero. This proves surjectivity. Strict convexity on $H$ proves injectivity.
\end{proof}
